\begin{abstract}
M1 established a proof--audit--witness contract for trustworthy finite impossibility evidence in the Sen24 base case; M1.5 addresses a different question, namely whether such evidence also determines a canonical repair presentation.
Working strictly in the Sen24 base case \((n,m)=(2,4)\) and within the current local-rationality setting induced by \code{no\_cycle3} and \code{no\_cycle4}, we compare a bundled liberalism lever \code{minlib} with a split encoding exposing \code{decisive\_voter0} and \code{decisive\_voter1}.
The repository artifact bundle shows that the bundled and split witness CNFs are clause-multiset equivalent under the mapped comparison and that both witnesses are UNSAT.
Nevertheless, their raw lever-level minimal repair families differ: the bundled witness exposes the singleton repair \code{minlib}, whereas the split witness exposes person-specific singleton repairs.
We formalize this as an existence theorem showing that raw repair explanation is not invariant under clause-equivalent encodings in the Sen24 base case.
The implication is normative rather than merely notational: finite witness legitimization can certify that a boundary is real, but it does not automatically determine how repairs should be grouped or presented.
Audit-oriented repair reporting therefore needs an explicit abstraction contract in addition to proof-carrying or audit-carrying witness infrastructure.
\end{abstract}
